\section{Moments}
\label{bdc:sec:moments}

We now turn from the fates of the cell to the numbers it carries and releases.
At first sight their moments appear awkward. Since the catastrophe rate grows
with the population, the equation for each load moment calls for the moment
above it. One would normally expect an open hierarchy and, sooner or later, a
closure approximation. The Riccati equation offers a way around this.

\begin{remark}[Why closed forms exist]
\label{bdc:rem:closure}
Let $J(t)=\ex{X_t}$ and $K(t)=\ex{X_t^2}$. Their moment equations begin
\[
\dot J=(\lambda-\mu)J-\delta K,
\qquad
\dot K=2(\lambda-\mu)K+(\lambda+\mu)J-\delta\ex{X_t^3},
\]
and the pattern continues. The hierarchy is genuinely open. However, the
non-catastrophe probability obeys the autonomous equation
\cref{bdc:eq:riccatifact}, while $J=-I'/\delta$. Differentiating that scalar
equation produces $K$ in terms of $I$ and $J$, without introducing the third
moment. Repeated differentiation does the same at higher orders. Thus the open
hierarchy remains present, but it need not be truncated; it can be approached
from the closed equation for $I$. \Cref{bdc:prop:polynomial} later states the
general result.
\end{remark}

The calculation has a natural order. We first recover the mean load from $I$,
then use its derivative to obtain the second moment. The release moments follow
from balance equations in which the same higher-order terms appear and can be
eliminated.

\subsection{The mean load}
\label{bdc:sec:J}

The link between $I$ and the load is especially direct. Since $H$ is entered
from state $n$ at rate $\delta n$, the forward equation for the catastrophe
state is
\[
\ddt{}\pr{X_t=H}
=\delta\sum_{n\ge1}n\,p_n(t)
=\delta J(t).
\]
Because $I(t)=1-\pr{X_t=H}$,
\begin{equation}
\label{bdc:eq:IJ}
\ddt{I}=-\delta J.
\end{equation}
The slope of the non-catastrophe curve is therefore the negative mean load,
scaled by the per-unit catastrophe rate. Differentiating \cref{bdc:eq:I}, or
substituting the Riccati equation into \cref{bdc:eq:IJ}, gives the closed form.

\begin{theorem}[The mean load]
\label{bdc:thm:J}
The expected number of replicative units present at time $t$ is
\begin{equation}
\label{bdc:eq:J}
J(t)=\ex{X_t}
=\frac{(a-b)^2\,w}{(B+Aw)^2}
=-\frac{\lambda}{\delta}(I-a)(I-b).
\end{equation}
\end{theorem}

The second expression will be the more useful one: it places $J$ directly on
the Riccati phase line. From the initial load of one, the mean may first rise as
productive paths grow, but it eventually returns to zero as absorption becomes
certain. This fall does not mean that the remaining productive cells are
empty. It means that there are progressively fewer of them, as
\cref{bdc:rem:Jmixes} made clear.

To expose the second moment, consider the expected change of $X_t$ over
$\Delta t$. A birth adds one at rate $\lambda X_t$, a death removes one at rate
$\mu X_t$, and a catastrophe removes the whole numerical load at rate
$\delta X_t$. Hence
\[
\ex{X_{t+\Delta t}}
=\ex{X_t}
+\ex{\lambda\Delta t\,X_t-\mu\Delta t\,X_t
-\delta\Delta t\,X_t^2}
+o(\Delta t),
\]
and therefore
\begin{equation}
\label{bdc:eq:ddJ}
\ddt{J}=(\lambda-\mu)J-\delta K(t),
\qquad K(t):=\ex{X_t^2}.
\end{equation}
The term $-\delta K$ is the distinctive one. A large load is both costly when
removed and more likely to be removed, so its contribution is quadratic.

\subsection{The second moment}
\label{bdc:sec:K}

Although \cref{bdc:eq:ddJ} couples $J$ to $K$, it can be read in reverse because
$J$ is already known. Differentiate \cref{bdc:eq:IJ}:
\begin{equation}
\label{bdc:eq:Kder1}
\ddt{J}=-\frac{1}{\delta}\frac{\dd^2I}{\dd t^2}.
\end{equation}
Differentiating the Riccati equation gives
\[
\frac{\dd^2I}{\dd t^2}
=\bigl[-(\lambda+\mu+\delta)+2\lambda I\bigr]\ddt{I},
\]
and using $I'=-\delta J$ once more,
\begin{equation}
\label{bdc:eq:Kder2}
\frac{\dd^2I}{\dd t^2}
=(\lambda+\mu+\delta)\delta J-2\lambda\delta I J.
\end{equation}
Equations~\cref{bdc:eq:Kder1,bdc:eq:Kder2} imply
\[
\ddt{J}=-(\lambda+\mu+\delta)J+2\lambda I J.
\]
Comparison with \cref{bdc:eq:ddJ} then gives
\[
(\lambda-\mu)J-\delta K
=-(\lambda+\mu+\delta)J+2\lambda I J,
\]
from which $K$ follows without ever computing $\ex{X_t^3}$.

\begin{theorem}[The second moment of the load]
\label{bdc:thm:K}
Let $\kappa=1+\delta/(2\lambda)$. Then
\begin{equation}
\label{bdc:eq:K}
K(t)=\ex{X_t^2}
=\Bigl[1+\frac{2\lambda}{\delta}(1-I)\Bigr]J
=\frac{2\lambda}{\delta}(\kappa-I)J,
\end{equation}
and, after substituting \cref{bdc:eq:J},
\begin{equation}
\label{bdc:eq:KI}
K(I)
=\frac{2\lambda^2}{\delta^2}(I-\kappa)(I-a)(I-b).
\end{equation}
\end{theorem}

As a polynomial in $I$, the second moment shares the roots $a$ and $b$ with
the mean. Its third root $\kappa$ lies just above unity, outside the interval
$(b,1]$ traced by the process.

\subsection{The mean release}
\label{bdc:sec:V}

The released count $W_t$ remains at zero until rupture and then takes the load
present at that instant. A rupture from state $n$ increases $W_t$ by $n$ at rate
$\delta n$, and so
\begin{equation}
\label{bdc:eq:ddV}
\ddt{V}=\delta K(t),
\qquad V(t):=\ex{W_t}.
\end{equation}
Release is driven by the second moment rather than the first: fuller cells
release more and are also more liable to rupture.

Combining \cref{bdc:eq:ddV} with \cref{bdc:eq:ddJ} eliminates $K$:
\[
\ddt{J}=(\lambda-\mu)J-\ddt{V}.
\]
Since $J=-I'/\delta$,
\[
\ddt{V}
=-\frac{\lambda-\mu}{\delta}\ddt{I}-\ddt{J},
\]
which can be integrated immediately.

\begin{theorem}[The mean release]
\label{bdc:thm:V}
With $V(0)=0$ and $J(0)=I(0)=1$,
\begin{equation}
\label{bdc:eq:V}
V(t)=\ex{W_t}
=\frac{\lambda-\mu}{\delta}[1-I(t)]-J(t)+1
=(1-I)\Bigl[1+\frac{\lambda}{\delta}(1-I)\Bigr].
\end{equation}
Its late-time value is
\begin{equation}
\label{bdc:eq:Vinf}
V_\infty
=\lim_{t\to\infty}V(t)
=\frac{\lambda-\mu}{\delta}(1-b)+1
=\frac{aB}{A}.
\end{equation}
When $\mu=0$, this reduces to $V_\infty=1+\lambda/\delta$.
\end{theorem}

The form in $I$ follows by substituting \cref{bdc:eq:J} into the first
expression and using \cref{bdc:eq:identities}. Notice that $V_\infty$ averages
over every realisation. Paths which die internally contribute zero, so it is not
yet the mean conditional burst size; that distinction is taken up in
\cref{bdc:sec:conditioning}.

\subsection{Load and release together}
\label{bdc:sec:JV}

\Cref{bdc:fig:jv} restores the two familiar time courses alongside the
parametric view. In time, $J$ rises and falls while $V$ accumulates. Eliminating
time shows the same movement as one trajectory in the $(J,V)$ plane. The two
panels are complementary: the first makes the timing legible, and the second
shows how intracellular load is exchanged for expected release.

\begin{figure}[!htbp]
\centering
\includegraphics[width=0.94\textwidth]{F4_JV_parametric.pdf}
\caption{Mean intracellular load and release for $\lambda=1$, $\mu=0.2$ and
$\delta=0.05$. (a) The time courses $J(t)$ and $V(t)$: the load first grows and
then disappears, while release increases to $V_\infty=13.99$. (b) The same
solution as the parametric path $(J(t),V(t))$, with selected times marked.}
\label{bdc:fig:jv}
\end{figure}

The second moment of the release is obtained by the same elimination: a
rupture from state $n$ raises $W_t^2$ from zero to $n^2$ at rate $\delta n$,
so $\ex{W_t^2}$ is driven by the third load moment, which is already present
in the equation for $K$. The calculation is in \cref{bdc:app:Wsq}.

\begin{theorem}[The second moment of the release]
\label{bdc:thm:Wsq}
The second moment of the released count is
\begin{equation}
\label{bdc:eq:Wsq}
\ex{W_t^2}
=\frac{2(\lambda-\mu)}{\delta}V
-\frac{\lambda+\mu}{\delta}I
-K+\frac{\lambda+\mu}{\delta}+1.
\end{equation}
\end{theorem}

\subsection{Variances}
\label{bdc:sec:variances}

The variances now require only
$\var{Y}=\ex{Y^2}-\ex{Y}^2$. For the load,
\begin{equation}
\label{bdc:eq:VarX}
\var{X_t}
=K(t)-J(t)^2
=\frac{2\lambda}{\delta}(\kappa-I)J-J^2,
\end{equation}
and for the release,
\begin{equation}
\label{bdc:eq:VarW}
\var{W_t}
=\frac{2(\lambda-\mu)}{\delta}V
-\frac{\lambda+\mu}{\delta}I
-K+\frac{\lambda+\mu}{\delta}+1-V^2.
\end{equation}
Since $J$, $K$ and $V$ are functions of $I$, the two variances are as well.

\begin{figure}[!htbp]
\centering
\includegraphics[width=0.94\textwidth]{F4_mean_sd_XW.pdf}
\caption{Means and one-standard-deviation envelopes for (a) the intracellular
load $X_t$ and (b) the released count $W_t$, with
$\lambda=1$, $\mu=0.2$ and $\delta=0.05$. Lower envelopes are clipped at the
support boundary zero.}
\label{bdc:fig:meansd}
\end{figure}

\Cref{bdc:fig:meansd} shows why these expressions are worth carrying. The
spread of the intracellular load soon exceeds its mean. At late times, variation
in $W_t$ reflects a mixture of paths that never ruptured and paths that released
a sometimes large load. Neither panel is well described by its central curve
alone.

\FloatBarrier

\subsection{The pair \texorpdfstring{$(X_t,W_t)$}{(Xt,Wt)}}
\label{bdc:sec:pair}

The two components have now been treated separately. Their joint behaviour has
a useful degeneracy: internal load and released load cannot be positive at the
same time.

\begin{corollary}[Exclusion, and the covariance]
\label{bdc:cor:exclusion}
At every $t$, $X_tW_t\equiv0$ almost surely under the numerical-load convention
of \cref{bdc:sec:notation}. Before rupture $W_t=0$, and after rupture the state
is $H$ with numerical load zero. Hence $\ex{X_tW_t}=0$ and
\begin{equation}
\label{bdc:eq:cov}
\cova{X_t,W_t}=-J(t)V(t)\le0.
\end{equation}
\end{corollary}

The sum $X_t+W_t$ has a similarly plain interpretation. It counts particles
currently held inside the intact cell together with particles already released.
It does not count particles lost through internal death, and it is not a
cumulative count of all births.

\begin{corollary}[Current internal and released yield]
\label{bdc:cor:yield}
The mean number of particles either still inside or already released is
\begin{equation}
\label{bdc:eq:yield}
\ex{X_t+W_t}
=1+\frac{\lambda-\mu}{\delta}[1-I(t)].
\end{equation}
\end{corollary}

Algebraically, \cref{bdc:eq:yield} is a rearrangement of
\cref{bdc:eq:V}. Probabilistically, it describes the surviving inventory of the
process: the founder, together with the net contribution of births and deaths
accumulated before time $t$ or rupture. This is the quantity whose late-time
limit underlies $V_\infty$.

\FloatBarrier
